% !TEX program = xelatex
\documentclass[12pt,fontset=none]{ctexart}
\usepackage[a4paper,margin=1in]{geometry}
\usepackage{amsmath,amssymb,amsthm,mathtools,bm}
\usepackage{booktabs,longtable,array}
\usepackage{enumitem}
\usepackage{setspace}
\usepackage{iftex}
\usepackage[colorlinks=true,linkcolor=blue,citecolor=blue,urlcolor=blue]{hyperref}

\setstretch{1.2}

\ifXeTeX\else\PackageError{coauthor-note}{Please compile this file with XeLaTeX to use FangSong Chinese fonts}{}\fi

\IfFontExistsTF{Noto Serif CJK SC}{
    \setCJKmainfont{Noto Serif CJK SC}[
        UprightFont=*,
        BoldFont=* Bold,
        ItalicFont=FandolKai-Regular
    ]
}{
    \setCJKmainfont{FandolSong}[
        UprightFont=*-Regular,
        BoldFont=*-Bold,
        ItalicFont=FandolKai-Regular
    ]
}
\IfFontExistsTF{Noto Sans CJK SC}{
    \setCJKsansfont{Noto Sans CJK SC}[
        UprightFont=*,
        BoldFont=* Bold
    ]
}{
    \setCJKsansfont{FandolHei}[
        UprightFont=*-Regular,
        BoldFont=*-Bold
    ]
}
\setCJKmonofont{FandolSong}[
    UprightFont=*-Regular,
    BoldFont=*-Bold
]

\title{给合作者的说明：\\
为什么我们要从 \texttt{A/B/C} 企业状态转移，转向“产品层面的 commercialization assignment”，\\
以及为什么政策进入模型的核心原语应改写为 \texorpdfstring{$\tau_{ext}$}{tau\_ext}
}
\author{}
\date{}

\begin{document}
\maketitle

\section*{说明目的}

这份说明用于解释：为什么我们现在需要把论文的核心解释对象，从
“企业身份如何在 \texttt{A/B/C} 之间变化”，更新为“一个原研药
idea / 产品最终由谁来实现”，以及为什么与这一调整相对应，
政策进入模型的方式也必须更新。

这次更新不是小修小补，而是对论文对象、识别重心和模型原语的一次统一调整。
最核心的新判断是：

\begin{quote}
\textbf{MAH 的作用不应再主要写成一个泛泛的制度 dummy 或 regime index，
而应写成 external commercialization 这条组织路线的
route-specific governance/compliance cost 的下降。}
\end{quote}

我们用
\[
\tau_{ext}
\]
来表示这个对象。它的含义不是“所有治理成本”，更不是“所有政策效应的总和”，
而是：

\begin{quote}
\textbf{研发型主体在保留 authorization / holder 身份的同时，
通过外部合格生产来完成 original-drug commercialization 所需要承担的额外治理与合规摩擦。}
\end{quote}

因此，这次调整实际上是两步同时完成：

\begin{enumerate}[leftmargin=2em]
    \item 把论文对象从 \texttt{firm-state-centered} 改成
    \texttt{product-level assignment-centered}；
    \item 把政策进入模型的方式，从“泛泛的制度变化”改成
    “external commercialization route-specific governance friction 的变化”。
\end{enumerate}

\section{一、旧设定在解释什么}

旧设定的主线是：

\begin{quote}
\texttt{MAH} 改变企业组织边界，企业在 \texttt{A/B/C} 三种组织状态之间重新选择，
进而影响原研药创新与行业均衡。
\end{quote}

这里的三类企业是：

\begin{itemize}[leftmargin=2em]
    \item \texttt{A}：研发与生产一体化企业；
    \item \texttt{B}：研发专门化企业；
    \item \texttt{C}：生产专门化企业。
\end{itemize}

在这个框架里，\texttt{MAH} 主要通过两个 reduced-form 渠道进入模型：

\begin{itemize}[leftmargin=2em]
    \item 降低 \texttt{B} 型企业的 commercialization friction，原来通常记为
    \(\tau_M\)；
    \item 降低企业在 \texttt{A/B/C} 之间切换的组织重构成本，记为
    \(\kappa_{sj}\)。
\end{itemize}

因此，旧设定最强调的经验对象是：

\begin{itemize}[leftmargin=2em]
    \item incumbent transition matrix \(P(M)\)；
    \item entry composition \(e(M)\)；
    \item 类型份额、退出、producer-side outcomes、commercialization proxies 等。
\end{itemize}

换句话说，旧版的识别逻辑更接近：

\begin{quote}
先看企业在 \texttt{A/B/C} 之间有没有明显重组，再把这种组织重组和创新结果联系起来。
\end{quote}

\section{二、新设定在解释什么}

新设定不再把“企业身份是否从 \texttt{A} 变成 \texttt{B/C}”当成主调整边际，
而是改成：

\begin{quote}
\texttt{MAH} 改变的是
\textbf{original-drug idea 从 invention 到 commercialization 的组织分配方式}。
\end{quote}

更直白地说，新设定关心的是：

\begin{itemize}[leftmargin=2em]
    \item 谁生成 original-drug idea；
    \item 谁持有 / 保留上市许可或产品权益；
    \item 谁来生产；
    \item 是内部一体化完成，还是外包给有资质的生产方，还是转给别的企业实现。
\end{itemize}

因此，新设定最核心的经济问题不再是：

\begin{quote}
\texttt{A/B/C} 企业身份有没有大规模跳转？
\end{quote}

而是：

\begin{quote}
一个 original-drug idea 最终是由谁、通过什么组织路线被实现出来？
\end{quote}

这和我们最近讨论的方向一致：现实中更容易观察到的，不是企业大规模从
\texttt{A/B/C} 改头换面，而是同一类企业围绕不同产品采取不同的
commercialization arrangement。

\section{三、为什么 policy enters the model 也必须跟着更新}

如果论文对象从“企业身份跃迁”改成“产品实现关系”，
那么政策进入模型的原语也必须跟着改变。

旧写法里，政策更容易被处理成：

\begin{itemize}[leftmargin=2em]
    \item 一个制度 dummy：\(M\in\{0,1\}\)；
    \item 或一个比较宽泛的制度 index；
    \item 或者同时进入很多 reduced-form 对象：\(\tau_M,\kappa_{sj},\phi_B,\ldots\)
\end{itemize}

这种写法在最初阶段有用，但现在的问题是：它太宽。
如果论文的核心机制已经被压缩到
“research-originated idea 如何跨企业边界完成 commercialization”，
那么 policy enters the model 的单一原语也应该对应到这条路线上。

因此，新的写法应该是：

\[
\tau_{ext}
\]

其中 \(\tau_{ext}\) 表示：

\begin{quote}
\textbf{research-originated original-drug idea 通过外部合格生产完成 commercialization 时，
该 route 特有的治理/合规摩擦。}
\end{quote}

这个对象包括：

\begin{itemize}[leftmargin=2em]
    \item holder 和 manufacturer 之间签约与执行的成本；
    \item 质量协议、委托协议、监督与协调成本；
    \item MAH 承担质量责任和上市放行责任所带来的组织成本；
    \item 跨企业完成 commercialization 的制度性执行成本。
\end{itemize}

它不包括：

\begin{itemize}[leftmargin=2em]
    \item 物理生产成本本身；
    \item 原研 idea generation 的研发成本；
    \item 行业需求或 market size；
    \item 所有企业共同面对的固定经营成本。
\end{itemize}

所以更准确地说：

\[
\tau_{ext}
=
\text{governance/compliance cost of the external commercialization route}
\]

而不是：

\[
\tau_{ext}
=
\text{all governance cost in the economy.}
\]

\section{四、为什么 \texorpdfstring{$\tau_{ext}$}{tau\_ext} 最贴 MAH 的制度现实}

这一步是整个调整里最关键的地方。

MAH 的功劳，不是让企业突然更会研发，也不是让药品需求突然变大；
它的功劳是把：

\begin{quote}
“没有厂的研发者想把药做出来”
\end{quote}

这件事，从必须依赖企业内部化解决，改成可以通过法律认可、责任清楚、
执行标准化的合格合同关系来解决。

换句话说，改革前，一个 research-originating firm 如果没有内部生产能力，
往往只能：

\begin{itemize}[leftmargin=2em]
    \item 自己往下游整合，去建厂、拿生产能力；
    \item 或者把项目 / 产品让给一体化企业去实现。
\end{itemize}

改革后，新的组织路线变得可行：

\begin{itemize}[leftmargin=2em]
    \item 研发者可以保留 holder / authorization；
    \item 通过外部合格生产企业完成制造；
    \item 同时由法律和监管规则明确质量责任、上市放行、质量协议和委托关系。
\end{itemize}

因此，MAH 的经济学内容最适合被解释为：

\begin{quote}
\textbf{它降低了跨企业完成 commercialization 的治理成本。}
\end{quote}

也就是：

\[
\tau_{ext}^{post}<\tau_{ext}^{pre}.
\]

\section{五、旧设定和新设定最本质的不同}

\subsection*{1. 主要调整单位不同}

旧设定的主要调整单位是\textbf{企业组织状态}。
它关心的是企业从 \texttt{A/B/C} 之间怎么切换。

新设定的主要调整单位是\textbf{产品 / idea 的实现方式}。
它关心的是一个 original-drug idea 最终：

\begin{itemize}[leftmargin=2em]
    \item 由 \texttt{A} 内部做出来；
    \item 由 \texttt{B} 保留权利、外包给 \texttt{A/C} 生产；
    \item 由 \texttt{B} 转给 \texttt{A} 实现；
    \item 或者根本没有被实现。
\end{itemize}

所以：

\begin{itemize}[leftmargin=2em]
    \item 旧设定是 \texttt{firm-state transition-centered}；
    \item 新设定是 \texttt{product-level commercialization assignment-centered}。
\end{itemize}

\subsection*{2. \texttt{A/B/C} 的含义不同}

旧设定里，\texttt{A/B/C} 更像是\textbf{动态状态变量}；
重点是企业是否从一个状态转到另一个状态。

新设定里，\texttt{A/B/C} 更像是\textbf{组织角色 / 能力载体}。
它们仍然重要，但它们不再必须表现为大规模身份跃迁，才能解释 MAH 的作用。

换句话说：

\begin{itemize}[leftmargin=2em]
    \item 旧版：\texttt{A/B/C} 是“谁变成谁”；
    \item 新版：\texttt{A/B/C} 是“谁在产品实现链条里扮演什么角色”。
\end{itemize}

\subsection*{3. 经验识别的重心不同}

旧设定最依赖的是：

\begin{itemize}[leftmargin=2em]
    \item pre/post transition matrix；
    \item incumbent switching；
    \item firm-level reorganization evidence。
\end{itemize}

新设定更依赖的是：

\begin{itemize}[leftmargin=2em]
    \item 新进入者是否更偏 \texttt{B}；
    \item original-drug products 是否更常由“研发方持证 + 外部合格生产”来实现；
    \item entrusted production / commissioned manufacturing / 持证与生产分离的文本和匹配证据；
    \item producer-side outcomes；
    \item new original-drug products 的实现数量。
\end{itemize}

因此，新设定不是不要 transition evidence，而是把它\textbf{降级为辅助证据}。
真正的一线证据，变成了\textbf{产品实现关系本身}。

\subsection*{4. 对 MAH 作用的理解不同}

旧设定更容易写成：

\begin{quote}
\texttt{MAH} 降低组织切换成本，因此企业从一体化走向专业化。
\end{quote}

新设定更准确地写成：

\begin{quote}
\texttt{MAH} 降低 original-drug idea 通过外部合格生产完成 commercialization 的治理摩擦，
使研发专门化主体即使不自建全部产能，也更有可能把 idea 变成产品。
\end{quote}

这个表述更贴近制度现实，也更贴近当前的数据条件。
因为我们现在更容易观察到“谁研发、谁持证、谁生产、谁实现产品”，
而不是企业身份的大规模跳转。

\section{六、为什么必须改，而不是继续坚持旧设定}

原因很简单：

\begin{quote}
\textbf{旧设定最核心的 empirical anchor——\texttt{A/B/C} 身份跃迁——
现在看起来并不是最强、最可观测、最可校准的 MAH 作用边际。}
\end{quote}

现实中更可能发生的是：

\begin{itemize}[leftmargin=2em]
    \item 新进入者更偏 \texttt{B}；
    \item \texttt{A} 继续做研发，但某些产品外包给有资质的生产方；
    \item \texttt{B} 做研发并保留某些权益，但找 \texttt{A/C} 生产；
    \item 还有一些产品转让、委托生产、合作 commercialization 的安排。
\end{itemize}

这些都说明：

\begin{quote}
MAH 更像是在\textbf{重组 original-drug idea 的实现路径}，
而不一定是在推动 incumbent 企业广泛改身份。
\end{quote}

所以，如果继续把 transition matrix 放在最中心的位置，就会出现一个问题：

\begin{itemize}[leftmargin=2em]
    \item 理论上很完整；
    \item 但经验上最难站住。
\end{itemize}

这就是为什么这次调整是必要的。

\section{七、哪些东西没有变}

这次调整并不是推翻原有模型，下面这些核心内容仍然保留：

\subsection*{1. \texttt{A/B/C} 三类组织角色仍然保留}

只是它们从“主要转移状态”更多变成“产品实现链中的组织角色”。

\subsection*{2. invention / implementation 分离仍然保留}

这一点其实比以前更重要。我们现在更明确地是在讲：

\begin{itemize}[leftmargin=2em]
    \item innovation 不是 idea generation 自动等于结果；
    \item 中间还有一个 implementation / commercialization 环节。
\end{itemize}

\subsection*{3. MAH 通过 commercialization friction 起作用仍然保留}

这一点反而比以前更核心。只是现在这层 friction 被更准确地写成：

\[
\tau_{ext}
\]

即 external commercialization route-specific governance/compliance cost。

\subsection*{4. 原研药创新仍然是最终结果变量}

主 outcome 仍然是 new original-drug products 及其相关 innovation outcomes。

\section{八、哪些东西实质性地变了}

\subsection*{1. 论文主问题变了}

旧版更像：

\begin{quote}
MAH 是否通过企业组织状态转移促进创新？
\end{quote}

新版更像：

\begin{quote}
MAH 是否通过重组 original-drug idea 从 invention 到 commercialization 的组织分配方式，
促进新原研药产品实现？
\end{quote}

\subsection*{2. 模型最核心的“状态变化”变了}

旧版：\texttt{A/B/C} 之间的企业状态转移。

新版：idea / 产品在不同组织角色之间的实现分配。

\subsection*{3. 数据工程重点变了}

旧版重点：transition matrices。

新版重点：持证—生产—产品实现链接、委托生产文本、产品层 commercialization evidence。

\subsection*{4. 论文解释口径变了}

旧版更像 \texttt{firm reorganization}。

新版更像 \texttt{organization of innovation / commercialization assignment}。

\subsection*{5. policy enters the model 的写法变了}

旧版更容易写成：
\[
M \in \{0,1\}
\]
并让 \(M\) 去影响很多 reduced-form 对象。

新版更稳的写法是：

\begin{quote}
政策最核心地进入 external commercialization 这一步，
通过降低 external-route governance/compliance friction：
\[
\tau_{ext}^{post}<\tau_{ext}^{pre}.
\]
\end{quote}

然后，再通过这一条 route 的静态 payoff 和动态 continuation value，
影响 \(x_B\)、\(V_B\)、Entry\(_B\) 和 new original-drug products。

\section{九、可以给合作者的最简洁概括}

最简洁的说法是：

\begin{quote}
\textbf{我们不是放弃 \texttt{A/B/C} 框架，而是把 \texttt{A/B/C}
从“企业身份跃迁模型”改成“产品实现关系模型”中的组织角色。}
\end{quote}

或者再明确一点：

\begin{quote}
\textbf{旧设定把 MAH 的作用主要理解为企业在 \texttt{A/B/C} 之间重组；
新设定把 MAH 的作用主要理解为 original-drug ideas 在 \texttt{A/B/C}
之间的 commercialization assignment 发生变化。}
\end{quote}

并且，和这一点对应地：

\begin{quote}
\textbf{政策进入模型的核心原语，也应从泛泛的制度 dummy，更新为
external commercialization 这条路线的治理/合规摩擦
\(\tau_{ext}\)。}
\end{quote}

\section{十、建议发给合作者时最后加的一句判断}

可以直接补这一句：

\begin{quote}
这次调整的核心理由，不是理论偏好，而是识别与校准：
现有数据更容易支持“产品实现关系变化”，而不容易支持“大规模企业身份跃迁”。
因此，把 MAH 的核心机制重写为
product-level commercialization assignment，
并把政策进入模型的核心对象改写为
external commercialization route-specific governance friction
\(\tau_{ext}\)，
比继续坚持 transition-matrix-centered 的版本更稳，也更贴制度现实。
\end{quote}

\end{document}
